\documentclass[a4paper, serif, 10pt]{moderncv}

%% moderncv
% style and color
\moderncvstyle{banking}
\moderncvcolor{black}

%% page layout/margins
\usepackage[top=2.5cm, bottom=2.5cm, left=1.5cm, right=1.5cm]{geometry}

\nopagenumbers{}

%% Babel config for German language
% ref:
% - https://latex3.github.io/babel/guides/locale-german.html
\usepackage[ngerman]{babel}

%% fontspec
\usepackage{fontspec}

\IfFontExistsTF{Linux Libertine}{
  \setmainfont[Ligatures={TeX, Common}, Numbers=OldStyle]{Linux Libertine}
}{}
\IfFontExistsTF{Linux Libertine O}{
  \setmainfont[Ligatures={TeX, Common}, Numbers=OldStyle]{Linux Libertine O}
}{}

%% CTeX
% CJK support, used to show CJK characters in the resume
%
% - fontset=none: disable builtin fontset but instead set the CJK font manually
% - heading=false: disable ctex heading
% - punct=kaiming: use kaiming punctuations styles for CJK
% - scheme=plain: use plain scheme, do not override \normalsize font size
% - space=auto: space settings for CJK characters
%
% ref:
% - http://ctan.mirrorcatalogs.com/language/chinese/ctex/ctex.pdf
\usepackage[UTF8, heading=false, punct=kaiming, scheme=plain, space=auto]{ctex}

\IfFontExistsTF{Noto Serif CJK SC}{
  \setCJKmainfont{Noto Serif CJK SC}
}{}
\IfFontExistsTF{Noto Sans CJK SC}{
  \setCJKsansfont{Noto Sans CJK SC}
}{}

%% line spacing
\usepackage{setspace}
\setstretch{1.125}

%% URL styling - use normal text instead of monospace
\urlstyle{same}

% Auto-underline all links
% Defer until \begin{document} so hyperref is already loaded
\AtBeginDocument{%
  \let\oldhref\href
  \renewcommand{\href}[2]{\oldhref{#1}{\underline{#2}}}%
}

%% Patch \httplink and \httpslink to handle full URLs
% The original moderncv macros always prepend http:// or https:// to URLs.
% This causes issues when the URL already has a protocol (e.g., https://example.com)
% resulting in malformed URLs like https://https://example.com.
% This patch checks if the URL already contains "://" and uses it directly if so.
% Note: We use \str_set:Nx (x-expansion) to expand macros like \@homepage first.
\ExplSyntaxOn
\str_new:N \g_yamlresume_url_str

\RenewDocumentCommand{\httpslink}{O{}m}{
  \str_set:Nx \g_yamlresume_url_str {#2}
  \str_if_in:NnTF \g_yamlresume_url_str {://}
    { \url{#2} }
    { \url{https://#2} }
}

\RenewDocumentCommand{\httplink}{O{}m}{
  \str_set:Nx \g_yamlresume_url_str {#2}
  \str_if_in:NnTF \g_yamlresume_url_str {://}
    { \url{#2} }
    { \url{http://#2} }
}
\ExplSyntaxOff

\name{Julia Becker}{}
\title{Customer Success Manager mit Fokus auf Kundenbindung und SaaS-Wachstum}
\phone[mobile]{+49 30 12345678}
\email{julia.becker@example.de}
\homepage{https://www.juliabecker.de}

\address{Friedrichstraße 123, Berlin, Berlin, Deutschland, 10117}{}{}

\extrainfo{{\small \faLinkedin}\ \href{https://www.linkedin.com/in/julia-becker-csm}{@julia-becker-csm} {} {} {} • {} {} {} 
{\small \faGithub}\ \href{https://github.com/juliabecker}{@juliabecker}}

\begin{document}

\maketitle

\section{Basisinformationen}

\cvline{}{Customer Success Manager mit über sechs Jahren Erfahrung in der Betreuung von B2B-Kunden im SaaS-Umfeld. Ich verbinde Datenanalyse mit empathischer Kommunikation, um Kundenziele zu erreichen und nachhaltige Geschäftsbeziehungen aufzubauen.
%
\begin{itemize}
\item Verantwortung für ein Portfolio mit einem jährlichen Umsatzvolumen von über 4 Mio. EUR
\item Erfolgreiche Steigerung der Kundenbindung um 18 \% durch proaktives Onboarding und QBRs
\item Ausgeprägte Stärke in Eskalationsmanagement, Stakeholder-Kommunikation und vertraglichen Verhandlungen
\item Technische Affinität zu CRM-, Ticketing- und Customer-Success-Plattformen wie Salesforce, Gainsight und Zendesk
\end{itemize}}
\section{Ausbildung}

\cventry{Okt. 2014–Sept. 2016}
        {Master, Wirtschaftsingenieurwesen, Punkte: 1,7}
        {Technische Universität Berlin}
        {\url{https://www.tu.berlin}}
        {}
        {\begin{itemize}
\item Schwerpunkt auf digitalen Geschäftsmodellen und kundenorientierter Unternehmensführung
\item Masterarbeit über Erfolgsfaktoren von Customer-Success-Organisationen im B2B-SaaS
\end{itemize}
\textbf{Kurse}: Strategisches Management, Marketing und Vertrieb, Datenanalyse für Entscheider, Projektmanagement, Verhandlungsführung, Service Design, Business Process Management, Customer Relationship Management}

\cventry{Okt. 2011–Sept. 2014}
        {Bachelor, Betriebswirtschaftslehre, Punkte: 2,0}
        {Humboldt-Universität zu Berlin}
        {\url{https://www.hu-berlin.de}}
        {}
        {\begin{itemize}
\item Fundierte betriebswirtschaftliche Grundlagen mit Fokus auf Marketing und Organisation
\item Praxisprojekt zur Kundenbindung im Einzelhandel
\end{itemize}
\textbf{Kurse}: Allgemeine Betriebswirtschaftslehre, Mikroökonomie, Statistik, Marketing, Personal und Organisation, Wirtschaftsinformatik}
\section{Berufserfahrung}

\cventry{Jan. 2022–Heute}
        {Senior Customer Success Manager}
        {CloudBridge GmbH}
        {\url{https://www.cloudbridge.de}}
        {}
        {\begin{itemize}
\item Betreuung eines Portfolios von 25 strategischen Enterprise-Kunden mit einem ARR von 5,2 Mio. EUR
\item Entwicklung und Umsetzung von Customer-Success-Plänen zur Steigerung der Produktadoption und Kundenzufriedenheit (NPS +22)
\item Durchführung von vierteljährlichen Business Reviews (QBRs) mit C-Level-Stakeholdern
\item Identifikation von Upsell- und Cross-Sell-Potenzialen in enger Zusammenarbeit mit dem Vertrieb
\item Proaktives Churn-Management durch Risikoanalyse und Eskalationspläne
\end{itemize}
\textbf{Schlüsselwörter}: Kundenbindung, Onboarding, QBR, Upselling, Churn-Management, Salesforce, Gainsight}

\cventry{März 2019–Dez. 2021}
        {Customer Success Manager}
        {DataFlow AG}
        {\url{https://www.dataflow.de}}
        {}
        {\begin{itemize}
\item Verantwortung für 60 mittelständische Kunden im DACH-Raum
\item Aufbau eines standardisierten Onboarding-Prozesses, der die Time-to-Value um 30 \% verkürzte
\item Analyse von Nutzungsdaten zur Ableitung konkreter Handlungsempfehlungen
\item Enge Zusammenarbeit mit Product und Support zur Weitergabe von Kundenfeedback
\item Erstellung von Health-Score-Modellen zur Früherkennung von Abwanderungsrisiken
\end{itemize}
\textbf{Schlüsselwörter}: SaaS, Onboarding, Health Score, Datenanalyse, Zendesk}

\cventry{Juli 2017–Feb. 2019}
        {Junior Customer Success Manager}
        {Nordwind Solutions}
        {\url{https://www.nordwind-solutions.de}}
        {}
        {\begin{itemize}
\item First-Level-Support für 120 Geschäftskunden per Telefon, E-Mail und Chat
\item Dokumentation von Kundenanfragen und Pflege der Wissensdatenbank
\item Vorbereitung von Reports und Präsentationen für das Customer-Success-Team
\item Mitwirkung an Retentionskampagnen und Kundenzufriedenheitsumfragen
\end{itemize}
\textbf{Schlüsselwörter}: Kundenservice, CRM, Reporting, Retention, Kommunikation}
\section{Sprachen}

\cvline{Deutsch}{Muttersprache oder zweisprachig \hfill \textbf{Schlüsselwörter}: Muttersprache}
\cvline{Englisch}{Verhandlungssichere Kenntnisse \hfill \textbf{Schlüsselwörter}: TOEFL 110, Verhandlungssicher}
\cvline{Spanisch}{Gute Kenntnisse \hfill \textbf{Schlüsselwörter}: Grundkenntnisse}
\section{Fähigkeiten}

\cvline{Customer Success Management}{Experte \hfill \textbf{Schlüsselwörter}: Kundenbindung, Onboarding, QBR, Churn-Management, Upselling}
\cvline{CRM und Tools}{Erfahren \hfill \textbf{Schlüsselwörter}: Salesforce, Gainsight, Zendesk, Jira, Confluence, Tableau}
\cvline{Kommunikation und Führung}{Erfahren \hfill \textbf{Schlüsselwörter}: Stakeholder-Management, Eskalationsmanagement, Verhandlung, Präsentation, Mentoring}
\section{Auszeichnungen}

\cventry{Nov. 2023}
        {CloudBridge GmbH}
        {Customer-Success-Manager des Jahres}
        {}
        {}
        {Auszeichnung für herausragende Leistungen bei der Steigerung der Kundenbindung und der Entwicklung innovativer Success-Programme.}

\cventry{Dez. 2020}
        {DataFlow AG}
        {Beste Teamleistung}
        {}
        {}
        {Anerkennung für kollaborative Zusammenarbeit und Unterstützung kollegenübergreifender Projekte.}
\section{Zertifikate}

\cventry{Juni 2021}
        {SuccessCOACHING}
        {Certified Customer Success Manager (CCSM)}
        {\url{https://www.successcoaching.com/certification}}
        {}
        {}

\cventry{Feb. 2020}
        {Salesforce}
        {Salesforce Certified Administrator}
        {\url{https://trailhead.salesforce.com/credentials/administrator}}
        {}
        {}

\cventry{Sept. 2022}
        {Scrum.org}
        {Professional Scrum Master I}
        {\url{https://www.scrum.org/professional-scrum-master-i-certification}}
        {}
        {}
\section{Veröffentlichungen}

\cventry{März 2022}
        {Kundenbindung im digitalen Zeitalter - Strategien für SaaS-Unternehmen}
        {Journal für Customer Management}
        {\url{https://www.example.com/publikation/kundenbindung-saas}}
        {}
        {\begin{itemize}
\item Analyse proaktiver Customer-Success-Ansätze zur Reduzierung von Churn
\item Fallstudien aus dem B2B-SaaS-Markt
\item Handlungsempfehlungen für Onboarding und QBRs
\end{itemize}}
\section{Projekte}

\cventry{Jan. 2023–Juni 2023}
        {Neugestaltung eines Self-Service-Portals zur Verbesserung der Kundenerfahrung}
        {Kundenportal Relaunch}
        {\url{https://www.cloudbridge.de/kundenportal}}
        {}
        {\begin{itemize}
\item Leitung eines cross-funktionalen Teams aus Product, Design und Support
\item Durchführung von Kundeninterviews und Usability-Tests
\item Steigerung der Portalnutzung um 45 \% innerhalb von drei Monaten
\end{itemize}
\textbf{Schlüsselwörter}: Customer Experience, Self-Service, Produktmanagement}

\cventry{März 2020–Nov. 2020}
        {Entwicklung eines datenbasierten Frühwarnsystems zur Churn-Prävention}
        {Health-Score-Modell}
        {\url{https://www.dataflow.de/health-score}}
        {}
        {\begin{itemize}
\item Integration von Nutzungs-, Support- und Umfragedaten
\item Definition von Schwellenwerten und automatisierten Warnmeldungen
\item Reduzierung der Churn-Rate um 12 \% im ersten Jahr
\end{itemize}
\textbf{Schlüsselwörter}: Datenanalyse, Churn, Tableau}
\section{Interessen}

\cvline{Sport}{Laufen, Yoga, Radfahren}
\cvline{Musik}{Klavier, Konzerte}
\cvline{Ehrenamt}{Mentoring, Bildung}
\section{Ehrenamtliche Tätigkeiten}

\cventry{Sept. 2020–Heute}
        {Mentorin für Berufseinsteiger}
        {Berliner Mentoring Netzwerk}
        {\url{https://www.berliner-mentoring.de}}
        {}
        {\begin{itemize}
\item Unterstützung von Berufseinsteigern bei der Karriereplanung und Bewerbungsstrategie
\item Durchführung monatlicher Coaching-Gespräche und Workshops
\item Aufbau eines Netzwerks für junge Fachkräfte im Bereich Customer Success
\end{itemize}}
    
\end{document}